% appendices/appendix_labs_and_exercises.tex
\chapter{Lab 汇总与动手实验指南}
\label{app:labs}

\section{环境准备}

所有 Lab 均基于 Python 3.9+，以下依赖已经足够：

\begin{lstlisting}[language=bash]
pip install torch numpy matplotlib  # 核心
pip install higher                   # Lab 3 (MAML) 可选
\end{lstlisting}

\section{Lab 索引}

\begin{center}
\begin{tabular}{lll}
\toprule
Lab & 对应章节 & 核心概念 \\
\midrule
Lab 2 & 第 \ref{ch:why_test_time} 章 & 分布偏移 + TTT 自监督微调 \\
Lab 3 & 第 \ref{ch:bilevel} 章 & MAML 正弦波 \\
Lab 4 & 第 \ref{ch:fastweights} 章 & 外积联想记忆 \\
Lab 5 & 第 \ref{ch:attention} 章 & 线性注意力递推 vs 朴素版 \\
Lab 6 & 第\ref{ch:icl} 章 & ICL 与最小二乘对比 \\
Lab 7 & 第 \ref{ch:ttt_bridge} 章 & TTT-RNN 单元（NumPy） \\
Lab 8 & 第 \ref{ch:ttt_layers} 章 & TTT-Linear vs GRU \\
Lab 9 & 第 \ref{ch:titans} 章 & Needle-in-Haystack 混合记忆 \\
Lab 10 & 第 \ref{ch:unified} 章 & 在线 OLS 与 TTT-Linear 对比 \\
Lab 11 & 第 \ref{ch:nested} 章 & NL 对应表（开放性思考） \\
\bottomrule
\end{tabular}
\end{center}

\section{Lab 设计原则}

每个 Lab 的设计遵循以下约束：
\begin{enumerate}
  \item 代码不超过 200 行（含注释）
  \item 运行时间不超过 5 分钟（CPU）
  \item 有明确的"预期观察"——学生应该知道自己在验证什么
  \item 失败或非预期结果同样值得分析——Lab 不是黑盒验证工具
\end{enumerate}
